\chapter{Results and Discussion}
\label{ch:results}

All models were evaluated on the same held-out test set of 56{,}746 transactions containing
95 frauds, which took no part in training or threshold selection. The chapter first reports
the model comparison, then the controlled resampling study, then the statistical testing
that forms the core of the project, and closes with what the results do and do not
establish.

\section{Model Comparison}
\label{sec:comparison}

Table~\ref{tab:comparison} ranks every pipeline by PR-AUC. The trivial baseline that always
predicts ``legitimate'' is included to show why accuracy is not used: it would score 99.83
percent accuracy while catching nothing, and its PR-AUC of 0.0017 is simply the prevalence
of the positive class.

\begin{table}[htbp]
    \centering
    \caption{Model comparison.}
    \label{tab:comparison}
    \small
    \begin{tabular}{l r r r r r r r r}
\hline
\textbf{Model} & \textbf{PR-AUC} & \textbf{ROC-AUC} & \textbf{Prec.} & \textbf{Recall} & \textbf{F1} & \textbf{TP} & \textbf{FP} & \textbf{FN} \\
\hline
\textbf{LightGBM + weighted} & \textbf{0.828} & \textbf{0.967} & \textbf{0.972} & \textbf{0.737} & \textbf{0.838} & \textbf{70} & \textbf{2} & \textbf{25} \\
XGBoost + weighted & 0.816 & 0.967 & 0.973 & 0.747 & 0.845 & 71 & 2 & 24 \\
Forest + weighted & 0.813 & 0.940 & 0.957 & 0.705 & 0.812 & 67 & 3 & 28 \\
Forest + SMOTE & 0.809 & 0.960 & 0.857 & 0.758 & 0.804 & 72 & 12 & 23 \\
XGBoost + SMOTE & 0.807 & 0.971 & 0.985 & 0.705 & 0.822 & 67 & 1 & 28 \\
LightGBM + SMOTE & 0.774 & 0.968 & 0.931 & 0.705 & 0.802 & 67 & 5 & 28 \\
LogReg + SMOTE & 0.689 & 0.957 & 0.837 & 0.758 & 0.796 & 72 & 14 & 23 \\
LogReg + weighted & 0.688 & 0.964 & 0.847 & 0.758 & 0.800 & 72 & 13 & 23 \\
Baseline (always legit) & 0.002 & 0.500 & 0.000 & 0.000 & 0.000 & 0 & 0 & 95 \\
\hline
\end{tabular}
\end{table}

From the table we can see that LightGBM with class weighting leads at PR-AUC 0.828,
followed by XGBoost with weighting (0.816) and Random Forest with weighting (0.813).
Logistic Regression trails in both of its variants at roughly 0.689, which fits with the
fact that it cannot represent interactions between the PCA components. All counts in the
table are taken at each model's own tuned threshold, and the winning row is set in bold.

Two patterns are visible. First, the gradient-boosted and forest models form a tight
leading group spanning just 0.021 of PR-AUC. Second, class weighting slightly beat SMOTE
for the stronger learners. LightGBM in particular drops from 0.828 to 0.774 when SMOTE
replaces class weighting, which is the largest single effect in the table.
Section~\ref{sec:resampling} isolates that effect properly, and
Section~\ref{sec:significance} tests whether any of these gaps are real.
Figure~\ref{fig:comparison} shows the same ranking graphically.

\begin{figure}[htbp]
    \centering
    \includegraphics[width=0.78\textwidth]{08_model_comparison.png}
    \caption{PR-AUC ranking.}
    \label{fig:comparison}
\end{figure}

From the figure above we can see how tightly the leading group is clustered, with only the
two Logistic Regression bars standing clearly apart from the rest. The baseline is so close
to zero that it is not visible at this scale.

\section{Beyond Precision and Recall}

Precision, recall, and $F_1$ do not cover everything a fraud team needs to know.
Table~\ref{tab:extended} adds six further metrics, derived from the same confusion matrices
and taken at each model's own tuned threshold.

\begin{table}[htbp]
    \centering
    \caption{Extended metrics.}
    \label{tab:extended}
    \small
    \begin{tabular}{l r r r r r r}
\hline
\textbf{Model} & \textbf{MCC} & \textbf{Spec.} & \textbf{NPV} & \textbf{Bal. acc.} & \textbf{Youden $J$} & \textbf{Alerts/1k} \\
\hline
\textbf{LightGBM + weighted} & \textbf{0.846} & \textbf{0.99996} & \textbf{0.99956} & \textbf{0.868} & \textbf{0.737} & \textbf{1.269} \\
XGBoost + weighted & 0.852 & 0.99996 & 0.99958 & 0.874 & 0.747 & 1.286 \\
Forest + weighted & 0.821 & 0.99995 & 0.99951 & 0.853 & 0.705 & 1.234 \\
Forest + SMOTE & 0.806 & 0.99979 & 0.99959 & 0.879 & 0.758 & 1.480 \\
XGBoost + SMOTE & 0.833 & 0.99998 & 0.99951 & 0.853 & 0.705 & 1.198 \\
LightGBM + SMOTE & 0.810 & 0.99991 & 0.99951 & 0.853 & 0.705 & 1.269 \\
LogReg + SMOTE & 0.796 & 0.99975 & 0.99959 & 0.879 & 0.758 & 1.516 \\
LogReg + weighted & 0.801 & 0.99977 & 0.99959 & 0.879 & 0.758 & 1.498 \\
\hline
\end{tabular}
\end{table}

From the table we can see three things. Specificity and \gls{npv} tell us almost nothing
here, because every
model scores above 0.999 on both, and it is the true-negative count of about 56{,}650 that
puts them there. MCC, which uses all four cells, tracks PR-AUC closely and confirms the
same ordering. The most practical column is alerts per 1{,}000 transactions, which
translates the models into staffing terms: the champion raises 1.27 alerts per thousand,
while Logistic Regression with SMOTE raises 1.52. That is a 20 percent larger review queue
for noticeably worse ranking quality.

\section{Operating Point of the Selected Model}

The selected LightGBM model uses a tuned threshold of 0.636. On the test set this gives the
confusion matrix below.

\begin{center}
\begin{tabular}{lrr}
\hline
& \textbf{Predicted fraud} & \textbf{Predicted legitimate} \\
\hline
\textbf{Actual fraud}       & 70 (caught)      & 25 (missed) \\
\textbf{Actual legitimate}  & 2 (false alarm)  & 56{,}649 (correct) \\
\hline
\end{tabular}
\end{center}

From the table we can see that the model catches 70 of the 95 frauds while raising only two
false alarms, which is a precision of 0.972 and a recall of 0.737. Almost every alarm is a
genuine fraud,
but roughly one in four frauds still slips through. Whether that balance is acceptable
depends on the institution's priorities, which is exactly why the threshold is exposed as a
live control in the web application rather than fixed in code.

\section{ROC and Precision-Recall Curves}

Figure~\ref{fig:roccurves} makes the argument of Section~2.3 concrete by plotting the same
eight models under both kinds of curve.

\begin{figure}[htbp]
    \centering
    \includegraphics[width=0.95\textwidth]{14_roc_and_pr_curves.png}
    \caption{ROC and PR curves.}
    \label{fig:roccurves}
\end{figure}

From the figure above we can see that on the left every model's ROC curve hugs the top-left
corner and all eight look excellent, with ROC-AUC values between 0.94 and 0.97, while on
the right the same models separate clearly and the weaker ones are exposed. The dashed line
on the right is the random-classifier baseline at the 0.167 percent prevalence. Both panels
describe identical predictions, and only the choice of axes differs.

This is the clearest justification for the metric choice made in
Chapter~\ref{ch:methodology}, and it is why every ranking in this report uses PR-AUC rather
than ROC-AUC.

\section{Resampling Strategy Study}
\label{sec:resampling}

\subsection{Design}

Table~\ref{tab:comparison} compares SMOTE against class weighting across four algorithms,
but the algorithm changes at the same time as the strategy, so the two effects are mixed
together. A controlled experiment is needed to isolate the resampling decision on its own.

The design holds everything else fixed. The classifier is LightGBM with the tuned
hyperparameters of Table~\ref{tab:hyper}, and the split, the seed, and the
threshold-selection procedure are identical to Chapter~\ref{ch:methodology}. Only the
resampling step in the pipeline varies:

\begin{enumerate}
  \item Class weighting, with no resampler at all; the loss is re-weighted instead.
  \item Random undersampling, discarding majority rows until the classes balance.
  \item Random oversampling, duplicating minority rows until the classes balance.
  \item SMOTE, interpolating synthetic minority points until the classes balance.
\end{enumerate}

Every sampler sits inside the pipeline, so it touches only the training fold.
Table~\ref{tab:resampling} reports the outcome of the four runs.

\begin{table}[htbp]
    \centering
    \caption{Resampling comparison.}
    \label{tab:resampling}
    \small
    \begin{tabular}{l r r r r r r}
\hline
\textbf{Strategy} & \textbf{PR-AUC} & \textbf{ROC-AUC} & \textbf{Prec.} & \textbf{Recall} & \textbf{F1} & \textbf{Train rows} \\
\hline
\textbf{Class weighting} & \textbf{0.828} & \textbf{0.967} & \textbf{0.972} & \textbf{0.737} & \textbf{0.838} & \textbf{170{,}235} \\
Random oversampling & 0.815 & 0.965 & 0.901 & 0.768 & 0.830 & 339{,}902 \\
SMOTE & 0.774 & 0.968 & 0.931 & 0.705 & 0.802 & 339{,}902 \\
Random undersampling & 0.548 & 0.966 & 0.766 & 0.379 & 0.507 & 568 \\
\hline
\end{tabular}
\end{table}

From the table we can see a spread of 0.280 in PR-AUC between the best and the worst
strategy, even though the algorithm and its hyperparameters never changed. The last column,
``train rows'', is the size of the training set the classifier actually sees after
resampling, and it explains most of that spread. Figures \ref{fig:respr} and \ref{fig:rescm}
show the same four runs as curves and as confusion matrices.

\begin{figure}[htbp]
    \centering
    \includegraphics[width=0.7\textwidth]{09_resampling_pr_curves.png}
    \caption{Resampling PR curves.}
    \label{fig:respr}
\end{figure}

From the figure above we can see that class weighting and random oversampling trace almost
the same curve, while random undersampling deteriorates sharply once recall exceeds about
0.4.

\begin{figure}[htbp]
    \centering
    \includegraphics[width=0.82\textwidth]{10_resampling_confusion_matrices.png}
    \caption{Resampling confusion matrices.}
    \label{fig:rescm}
\end{figure}

From the figure above we can see the same result in counts: the three strategies that keep
all the data catch 67 to 73 of the 95 frauds, while undersampling catches only 36 of them.

\subsection{Interpretation}

Class weighting wins at 0.828. This confirms that its advantage in
Table~\ref{tab:comparison} was not caused by which algorithms happened to be paired with it.

Random oversampling is a close second at 0.815, only 0.013 behind. That makes
sense, because duplicating minority rows and re-weighting the loss are almost the same
operation. Duplicating a row $k$ times multiplies its contribution to the total loss by $k$,
which is exactly what a class weight of $k$ does directly, and without the memory cost of
building a table of 339{,}902 rows. Oversampling reaches the highest recall of the four
(0.768) at some cost in precision (0.901 against 0.972).

SMOTE underperforms both at 0.774. Its synthetic points are interpolated between
existing frauds, so they fill in regions the model already covers instead of extending what
it knows. There is a second problem. In a 29-dimensional space where the minority class is
sparse, interpolating between two frauds that are nearest neighbours but still far apart can
create points in regions where no real fraud lives, which adds label noise.

Random undersampling collapses to 0.548, and the reason is visible in the last
column of Table~\ref{tab:resampling}. Balancing 284 training frauds against 169{,}951
legitimate transactions means discarding 169{,}667 of the latter, which leaves a training
set of just 568 rows, or 0.33 percent of the data. Recall falls to 0.379, so it
misses 59 of 95 frauds, worse than every alternative. Its ROC-AUC nevertheless stays at
0.966, essentially tied with the winners. That is a clear demonstration that ROC-AUC can
hide a badly broken model on imbalanced data, exactly as Section~2.3 predicts.

For a dataset this imbalanced, the cost of throwing away majority data is greater than the
cost of the imbalance itself. The methods that keep all the data and adjust the learning
objective, namely class weighting and its near-equivalent oversampling, beat the methods
that change the data distribution.

\section{Confidence Intervals}
\label{sec:ci}

Sections \ref{sec:comparison} and \ref{sec:resampling} reported single numbers from a single
split. The rest of this chapter asks whether those numbers support the conclusions usually
drawn from them. The short answer is that they support much less than they appear to.

Table~\ref{tab:ci} attaches a 95 percent range to every headline figure, using a
1{,}000-resample percentile bootstrap for PR-AUC and Wilson score intervals for precision
and recall.

\begin{table}[htbp]
    \centering
    \caption{Confidence intervals.}
    \label{tab:ci}
    \small
    \begin{tabular}{l l l l}
\hline
\textbf{Model} & \textbf{PR-AUC [95\% CI]} & \textbf{Precision [95\% CI]} & \textbf{Recall [95\% CI]} \\
\hline
\textbf{LightGBM + weighted} & \textbf{0.828 [0.751, 0.899]} & \textbf{0.972 [0.904, 0.992]} & \textbf{0.737 [0.640, 0.815]} \\
XGBoost + weighted & 0.816 [0.738, 0.890] & 0.973 [0.905, 0.992] & 0.747 [0.652, 0.824] \\
Forest + weighted & 0.813 [0.731, 0.884] & 0.957 [0.881, 0.985] & 0.705 [0.607, 0.788] \\
Forest + SMOTE & 0.809 [0.734, 0.879] & 0.857 [0.767, 0.916] & 0.758 [0.663, 0.833] \\
XGBoost + SMOTE & 0.807 [0.733, 0.880] & 0.985 [0.921, 0.997] & 0.705 [0.607, 0.788] \\
LightGBM + SMOTE & 0.774 [0.688, 0.854] & 0.931 [0.848, 0.970] & 0.705 [0.607, 0.788] \\
LogReg + SMOTE & 0.689 [0.581, 0.789] & 0.837 [0.745, 0.900] & 0.758 [0.663, 0.833] \\
LogReg + weighted & 0.688 [0.580, 0.785] & 0.847 [0.756, 0.908] & 0.758 [0.663, 0.833] \\
\hline
\end{tabular}
\end{table}

From the table we can see that every interval is far wider than the gaps between the models
it is meant to separate. Figure~\ref{fig:forest} plots the PR-AUC intervals side by side.

\begin{figure}[htbp]
    \centering
    \includegraphics[width=0.78\textwidth]{11_pr_auc_confidence_intervals.png}
    \caption{PR-AUC confidence intervals.}
    \label{fig:forest}
\end{figure}

From the figure above we can see that the 95 percent bootstrap intervals of the top five
models overlap almost completely, so their bars give no visual evidence of a real ordering.

The intervals are wide. The champion's PR-AUC of 0.828 carries an interval of
$[0.751, 0.899]$, a span of 0.148. That is seven times the 0.021 that separates the entire
leading group of five models. Every one of those five has an interval that contains the
point estimates of the other four. Even Logistic Regression, which looks far behind at
0.689, has an upper bound of 0.789 that comfortably exceeds the champion's lower bound of
0.751.

The reason is the small number of positives. With 95 frauds in the test set, a single fraud
crossing the decision boundary changes recall by 1.05 percentage points. The recall
intervals are correspondingly wide, and the champion's spans 17 points, from 0.640 to
0.815.

\section{Pairwise Significance Testing}
\label{sec:significance}

Overlapping confidence intervals are suggestive but not decisive, because the models are
evaluated on the same transactions and their errors are correlated. McNemar's exact test is
the correct paired procedure (Section~\ref{sec:statmethods}). Table~\ref{tab:sig} shows the
ten pairs that came closest to significance out of the 28.

\begin{table}[htbp]
    \centering
    \caption{McNemar test results.}
    \label{tab:sig}
    \small
    \begin{tabular}{l l r r r}
\hline
\textbf{Model A} & \textbf{Model B} & \textbf{Discordant} & \textbf{$p$} & \textbf{$p_{\mathrm{Holm}}$} \\
\hline
XGBoost + weighted & LogReg + SMOTE & 13 & 0.0034 & 0.096 \\
LogReg + weighted & XGBoost + weighted & 12 & 0.0063 & 0.171 \\
LightGBM + weighted & LogReg + SMOTE & 14 & 0.0129 & 0.336 \\
LogReg + weighted & LightGBM + weighted & 13 & 0.0225 & 0.562 \\
XGBoost + weighted & Forest + SMOTE & 17 & 0.0490 & 1.000 \\
XGBoost + weighted & LightGBM + SMOTE & 13 & 0.0923 & 1.000 \\
LogReg + SMOTE & XGBoost + SMOTE & 18 & 0.0963 & 1.000 \\
LightGBM + weighted & Forest + SMOTE & 18 & 0.0963 & 1.000 \\
Forest + weighted & XGBoost + weighted & 7 & 0.1250 & 1.000 \\
LogReg + weighted & XGBoost + SMOTE & 17 & 0.1435 & 1.000 \\
\hline
\end{tabular}
\end{table}

From the table we can see that even the ten closest comparisons stay above the 0.05 level
once the $p$-values are Holm-adjusted. The full set of 28 comparisons is in
Appendix~\ref{app:tables}, and Figure~\ref{fig:mcnemar} shows all of them at once.

\begin{figure}[htbp]
    \centering
    \includegraphics[width=0.8\textwidth]{12_mcnemar_significance.png}
    \caption{McNemar significance matrix.}
    \label{fig:mcnemar}
\end{figure}

From the figure above we can see that no cell in the matrix falls below 0.05, so not one of
the 28 pairs differs significantly at the 5 percent level after Holm correction. The
smallest raw $p$-value is 0.0034, comparing XGBoost + weighted against
Logistic Regression + SMOTE, and it rises to 0.096 after correcting for 28 comparisons. The
comparison that actually matters, LightGBM + weighted against XGBoost + weighted, is
nowhere near significance.

The reason is visible in the discordant-pair counts: the models disagree on only 12 to 17
transactions out of 56{,}746. On that little evidence no test could tell them apart. This is
not a weakness of the test. It is an accurate statement about how much information 95
positive cases can carry.

\section{Cross-Validation}
\label{sec:cv}

A single split gives one sample of a quantity that varies. Five-fold stratified
cross-validation on the combined training and validation data re-estimates each model's
PR-AUC five times.

\begin{table}[htbp]
    \centering
    \caption{Cross-validation results.}
    \label{tab:cv}
    \small
    \begin{tabular}{l r r r r}
\hline
\textbf{Model} & \textbf{Mean PR-AUC} & \textbf{Std.} & \textbf{Min} & \textbf{Max} \\
\hline
\textbf{XGBoost + weighted} & \textbf{0.8494} & \textbf{0.0264} & \textbf{0.8182} & \textbf{0.8912} \\
LightGBM + weighted & 0.8490 & 0.0320 & 0.8138 & 0.8956 \\
Forest + SMOTE & 0.8471 & 0.0179 & 0.8212 & 0.8695 \\
Forest + weighted & 0.8441 & 0.0240 & 0.8039 & 0.8659 \\
\hline
\end{tabular}
\end{table}

From the table we can see that the four mean scores sit within 0.005 of each other while
each model's own standard deviation across folds is several times larger than that.
Figure~\ref{fig:cv} shows the same numbers as a distribution.

\begin{figure}[htbp]
    \centering
    \includegraphics[width=0.78\textwidth]{17_cross_validation.png}
    \caption{Cross-validation spread.}
    \label{fig:cv}
\end{figure}

From the figure above we can see that the spread within each model is wider than the
distance between the models, so the boxes overlap heavily.

Two results stand out.

First, the fold-to-fold standard deviation of 0.018 to 0.032 is four to six times the 0.005
that separates the four models' means. Any of these models can look best or worst depending
on which fold it is evaluated on. LightGBM ranges from 0.814 to 0.896 across folds, a span
wider than the entire distance between the best and worst model in the table.

Second, the ranking changes. XGBoost + weighted achieves the highest mean (0.8494)
against LightGBM's 0.8490, reversing the single-split result that selected LightGBM as
champion. The difference is 0.0004, far inside the noise, so the fair statement is not that
XGBoost is better but that the two cannot be told apart and the original choice was
arbitrary.

The cross-validated means (0.844 to 0.849) are higher than the single-split test scores
(0.813 to 0.828). That is because each fold trains on 80 percent of a larger pool than the
original 60 percent training split.

\section{Calibration}
\label{sec:calibration}

PR-AUC measures ranking. It says nothing about whether a predicted probability of 0.8 really
corresponds to an 80 percent chance of fraud, which matters as soon as the score feeds an
expected-loss calculation. The \gls{ece} \citep{guo2017calibration} checks this directly: it
groups predictions into probability bins and measures, on average, how far each bin's
predicted confidence is from the fraction of cases in that bin that turned out to be fraud.
A well-calibrated model has probabilities that can be trusted at face value, and a poorly
calibrated one does not, even if it ranks transactions correctly.
Table~\ref{tab:calibration} reports the ECE and the Brier score for all eight models,
ranked by ECE.

\begin{table}[htbp]
    \centering
    \caption{Calibration metrics.}
    \label{tab:calibration}
    \small
    \begin{tabular}{l r r}
\hline
\textbf{Model} & \textbf{ECE} & \textbf{Brier score} \\
\hline
\textbf{Forest + weighted} & \textbf{0.00034} & \textbf{0.00050} \\
XGBoost + weighted & 0.00042 & 0.00045 \\
LightGBM + weighted & 0.00046 & 0.00046 \\
LightGBM + SMOTE & 0.00065 & 0.00064 \\
XGBoost + SMOTE & 0.00106 & 0.00087 \\
Forest + SMOTE & 0.00252 & 0.00061 \\
LogReg + weighted & 0.06910 & 0.02374 \\
LogReg + SMOTE & 0.07102 & 0.02603 \\
\hline
\end{tabular}
\end{table}

From the table we can see that the six tree-based models occupy the top of the ranking and
the two Logistic Regression variants sit at the bottom, separated from them by two orders
of magnitude. Figure~\ref{fig:calibration} shows what that difference looks like.

\begin{figure}[htbp]
    \centering
    \includegraphics[width=0.72\textwidth]{13_calibration_curves.png}
    \caption{Reliability curves.}
    \label{fig:calibration}
\end{figure}

From the figure above we can see which models can be trusted at face value: points lying on
the diagonal are those whose predicted probabilities match the observed fraud frequencies,
and the curves that bend away from it are overconfident.

The tree-based models are extremely well calibrated, with ECE between 0.00034 and
0.00252 and Brier scores below 0.0009. Logistic Regression is not: its ECE of 0.069 is two
orders of magnitude worse, and its Brier score of 0.024 is 50 times the champion's.

This is a genuinely useful finding that the ranking metrics hide. Logistic Regression is
only moderately worse than the tree models at ordering transactions (PR-AUC 0.689 against
0.828), but it is far worse at quantifying risk. Its tuned threshold of 0.99999 in
Table~\ref{tab:comparison} is the symptom. The model is so consistently overconfident that
a sensible operating point exists only in the last fraction of a percent of its probability
range.

\section{Feature Importance}
\label{sec:importance}

Table~\ref{tab:importance} ranks the top twelve features by how much the model actually
relies on them, with its own gain-based importance alongside for comparison.

\begin{table}[htbp]
    \centering
    \caption{Feature importance.}
    \label{tab:importance}
    \small
    \begin{tabular}{l r r r r}
\hline
\textbf{Feature} & \textbf{Perm. importance} & \textbf{Std.} & \textbf{Gain \%} & \textbf{Splits} \\
\hline
V14 & 0.0726 & 0.0112 & 50.91 & 564 \\
V4 & 0.0541 & 0.0083 & 7.56 & 454 \\
V7 & 0.0140 & 0.0025 & 0.36 & 399 \\
V28 & 0.0120 & 0.0027 & 0.73 & 363 \\
V3 & 0.0117 & 0.0050 & 1.93 & 450 \\
V26 & 0.0112 & 0.0046 & 0.95 & 502 \\
V11 & 0.0104 & 0.0045 & 0.74 & 370 \\
V12 & 0.0087 & 0.0053 & 9.65 & 406 \\
Amount & 0.0082 & 0.0052 & 0.97 & 540 \\
V27 & 0.0061 & 0.0022 & 0.35 & 346 \\
V13 & 0.0051 & 0.0024 & 0.97 & 422 \\
V20 & 0.0049 & 0.0021 & 0.75 & 270 \\
\hline
\end{tabular}
\end{table}

From the table we can see that the two measures of importance broadly agree on which
features matter, but not on how much, and that both are dominated by a small number of
components. Figure~\ref{fig:importance} places them side by side.

\begin{figure}[htbp]
    \centering
    \includegraphics[width=0.95\textwidth]{15_feature_importance.png}
    \caption{Feature importance comparison.}
    \label{fig:importance}
\end{figure}

From the figure above we can see the drop-off in both panels: a couple of tall bars
followed by a long tail of features that contribute very little.

Permutation importance is the drop in test PR-AUC when a column is randomly shuffled. It is
the more trustworthy measure, because it is computed on data the model never saw. It
identifies V14 as dominant (0.0726) followed by V4 (0.0541), with a steep
drop to V7 (0.0140).

This disagrees with the exploratory analysis of Section~\ref{sec:eda}, which ranked V17
highest by linear correlation, with V14 second. The disagreement is informative. Univariate
correlation measures how strongly a feature is linearly associated with the label on its
own, while permutation importance measures what a feature adds given every other feature in
a non-linear model. V17 and V14 are both informative, but once V14 is present the model no
longer needs V17. Any feature-selection decision based on the correlation ranking alone
would have been misleading.

The gain column shows V14 accounting for 50.9 percent of total split gain, which confirms
that the model concentrates heavily on a small number of components.

\section{Learning Curve}

Table~\ref{tab:lc} retrains the champion on increasing fractions of the training set.

\begin{table}[htbp]
    \centering
    \caption{Learning curve results.}
    \label{tab:lc}
    \small
    \begin{tabular}{r r r r}
\hline
\textbf{Fraction} & \textbf{Training rows} & \textbf{Frauds} & \textbf{Test PR-AUC} \\
\hline
5\% & 8{,}511 & 15 & 0.736 \\
10\% & 17{,}023 & 33 & 0.740 \\
25\% & 42{,}558 & 75 & 0.803 \\
50\% & 85{,}117 & 159 & 0.802 \\
75\% & 127{,}676 & 227 & 0.812 \\
100\% & 170{,}235 & 284 & 0.828 \\
\hline
\end{tabular}
\end{table}

From the table we can see that PR-AUC keeps improving at every step, including the last one,
so the model has not yet run out of use for additional data. Figure~\ref{fig:lc} plots the
same values.

\begin{figure}[htbp]
    \centering
    \includegraphics[width=0.72\textwidth]{18_learning_curve.png}
    \caption{Learning curve.}
    \label{fig:lc}
\end{figure}

From the figure above we can see that the curve is still climbing at the right-hand edge
rather than levelling off, even on a logarithmic axis.

Performance rises from 0.736 at 5 percent of the training data (15 frauds) to 0.828 at 100
percent (284 frauds). Critically, the curve has not flattened, and the last
increment from 75 to 100 percent still adds 0.016. The binding constraint is the absolute
number of fraud examples, not the number of transactions. This makes collecting more
labelled fraud the single highest-value improvement available, ahead of any further
algorithmic tuning.

\section{Lift and Operational Value}

Table~\ref{tab:lift} sorts the test set by the champion model's predicted score and splits
it into ten equal groups.

\begin{table}[htbp]
    \centering
    \caption{Decile lift.}
    \label{tab:lift}
    \small
    \begin{tabular}{r r r r r}
\hline
\textbf{Decile} & \textbf{Reviewed} & \textbf{Frauds caught} & \textbf{Cum. recall} & \textbf{Lift} \\
\hline
1 & 5{,}674 & 88 & 0.926 & 9.26 \\
2 & 11{,}349 & 88 & 0.926 & 4.63 \\
3 & 17{,}023 & 89 & 0.937 & 3.12 \\
4 & 22{,}698 & 92 & 0.968 & 2.42 \\
5 & 28{,}373 & 94 & 0.989 & 1.98 \\
6 & 34{,}047 & 94 & 0.989 & 1.65 \\
7 & 39{,}722 & 95 & 1.000 & 1.43 \\
8 & 45{,}396 & 95 & 1.000 & 1.25 \\
9 & 51{,}071 & 95 & 1.000 & 1.11 \\
10 & 56{,}746 & 95 & 1.000 & 1.00 \\
\hline
\end{tabular}
\end{table}

From the table we can see that almost all of the fraud is concentrated in the first decile,
and that the deciles below it contain very little. Figure~\ref{fig:lift} shows the same
concentration two ways.

\begin{figure}[htbp]
    \centering
    \includegraphics[width=0.95\textwidth]{19_lift_and_gains.png}
    \caption{Cumulative gains and lift.}
    \label{fig:lift}
\end{figure}

From the figure above we can see that the gains curve on the left rises almost vertically
over the first decile and is then nearly flat, while the lift bars on the right drop away
sharply after the first group.

Framed as a review queue, the results are strong. The top decile of scored
transactions contains 88 of the 95 frauds, which is 92.6 percent of all fraud for 10
percent of the review effort, a lift of 9.26 over random sampling. The second decile adds
nothing at all, and the remaining seven frauds are scattered thinly through the rest of the
distribution.

This puts the recall figure in a useful light. At its tuned threshold the model reports
recall 0.737, but that reflects a very conservative cut-off chosen to hold false alarms near
zero. An institution willing to review 10 percent of transactions could reach 92.6 percent
recall with the same model and no retraining, purely by moving the threshold.

\section{Cost-Sensitive Thresholding}
\label{sec:cost}

The threshold objective of Section~\ref{sec:threshold} treats one missed fraud as equal to
one false alarm. A monetary model is more defensible: a missed fraud costs the transaction
amount actually lost, while a false alarm costs a fixed manual review, assumed here at
EUR 3. Let $\mathcal{M}(t)$ be the set of frauds the model misses at threshold $t$, let
$a_i$ be the amount of transaction $i$, and let $FP(t)$ be the number of false alarms. The
total cost is then
\begin{equation}
  \label{eq:cost}
  C(t) = \sum_{i \in \mathcal{M}(t)} a_i \; + \; 3 \, FP(t).
\end{equation}

Equation~\ref{eq:cost} is this project's instance of the expected-cost formulation of
\citet{elkan2001foundations}, with the general cost matrix given there filled in for this
problem: the cost of a missed fraud is the transaction amount and the cost of a false alarm
is the flat review fee. The EUR 3 figure is our own assumption rather than a published one.
The cost is evaluated across the full range of thresholds, and Table~\ref{tab:cost} reports
it at a set of representative ones.

\begin{table}[htbp]
    \centering
    \caption{Cost analysis.}
    \label{tab:cost}
    \small
    \begin{tabular}{r r r r r r}
\hline
\textbf{Threshold} & \textbf{Missed} & \textbf{False alarms} & \textbf{Fraud loss (EUR)} & \textbf{Review (EUR)} & \textbf{Total (EUR)} \\
\hline
0.005 & 20 & 9 & 3{,}786 & 27 & 3{,}813 \\
0.010 & 20 & 7 & 3{,}786 & 21 & 3{,}807 \\
0.050 & 22 & 5 & 3{,}808 & 15 & 3{,}823 \\
0.250 & 24 & 3 & 3{,}813 & 9 & 3{,}822 \\
0.500 & 25 & 2 & 3{,}813 & 6 & 3{,}819 \\
0.635 & 25 & 2 & 3{,}813 & 6 & 3{,}819 \\
0.750 & 25 & 2 & 3{,}813 & 6 & 3{,}819 \\
0.950 & 26 & 2 & 4{,}053 & 6 & 4{,}059 \\
\hline
\end{tabular}
\end{table}

From the table we can see that the cost-optimal point and the error-count threshold
actually deployed end up within a few euros of each other, despite sitting at very
different places on the threshold range. Figure~\ref{fig:cost} shows why.

\begin{figure}[htbp]
    \centering
    \includegraphics[width=0.82\textwidth]{16_cost_analysis.png}
    \caption{Cost against threshold.}
    \label{fig:cost}
\end{figure}

From the figure above we can see that the total-cost curve has a very shallow minimum: it
rises steeply only at the extreme low end, where every transaction is reviewed, and is
close to flat across the whole usable range of thresholds.

The cost-optimal threshold is 0.010, giving a total cost of EUR 3{,}807. The deployed
error-count threshold of 0.636 costs EUR 3{,}819, a difference of EUR 12, or 0.3
percent. The simple objective is, on this data, essentially cost-optimal.

The reason is that the total-cost curve is flat across most of its range, and the reason for
that is the finding from Section~\ref{sec:eda} that fraudulent transactions are typically
small, with a median of EUR 9.82. Because the frauds this model misses are mostly low in
value, the extra fraud loss from a stricter threshold is small, while review costs stay
negligible at EUR 3 apiece. Both endpoints confirm the shape: at $t = 0$ every transaction
is reviewed, costing EUR 169{,}953, and at $t = 1$ nothing is caught, costing EUR 14{,}766
in fraud losses.

This result depends on the assumed EUR 3 review cost. A higher figure would push the optimum
toward stricter thresholds, and the analysis should be repeated with a real cost.

\section{What the Results Establish}

Three findings in this project rest on evidence strong enough to state plainly.

Resampling strategy matters more than the choice of algorithm. The controlled
study of Section~\ref{sec:resampling} produced a spread of 0.280 PR-AUC between the best and
worst strategy with the algorithm held constant, while the entire leading group of
algorithms spans only 0.021. The reason is easy to see: random undersampling cuts the
training set from 170{,}235 rows to 568.

Tree ensembles are much better calibrated than logistic regression. An ECE gap of
two orders of magnitude, 0.0005 against 0.069, is not a marginal difference, and unlike the
ranking differences it is not drowned out by sampling noise.

The model is useful in practice whichever champion is chosen. Putting 92.6 percent
of frauds into the top-scoring decile is a large practical gain, and it holds for any of the
top models.

\section{What the Results Do Not Establish}
\label{sec:notestablish}

Our results do not show that LightGBM is better than XGBoost, Random Forest, or their SMOTE
variants.

Three separate checks point the same way. None of the 28 pairwise McNemar comparisons is
significant after Holm correction. The bootstrap intervals for the top five models overlap
almost completely. And when we move from a single split to five-fold cross-validation,
XGBoost comes out ahead instead.

The gap we would otherwise have reported is 0.012 of PR-AUC between LightGBM and XGBoost,
and those two models disagree on only 12 of the 56{,}746 test transactions. That is not
enough evidence to call a winner.

We think this is a useful result rather than a disappointing one. It tells us what 473 fraud
cases can and cannot support, and it is the kind of thing a single train/test split would
never have shown us.

\section{Threats to Validity}

\subsection{Internal Validity}

The hyperparameter tuning is uneven. Only LightGBM received a search of 40
candidates, while the other three algorithms use fixed defaults. The comparison is therefore
biased in LightGBM's favour, which makes its failure to separate from the others an even
weaker basis for choosing it.

The threshold is tuned on a small validation set, which contains only 94 frauds.
A threshold chosen to minimise total errors on 94 positives carries a lot of variance, and
its apparent optimality will not transfer exactly.

The test set is also reused. It was used once for the main comparison,
but the analyses from Section~\ref{sec:ci} onward reuse it for bootstrap intervals,
permutation importance, calibration, and cost analysis. These are descriptive analyses of a
fixed model rather than selection decisions, so they do not bias which model we chose. The
confidence intervals are nevertheless conditional on this particular test set, and are not
independent of one another.

\subsection{External Validity}

The data covers a single two-day window. All of it comes from September 2013 and from
one European card issuer. Fraud patterns drift, so nothing here can say anything about
performance a year later, in another market, or under a different fraud regime.

The features are anonymised. The PCA components cannot be interpreted, so the model's
reliance on V14 and V4 cannot be sanity-checked against domain knowledge. If those
components encode something specific to this collection window, the model would not
transfer, and there is no way to detect that from the data available.

There is no per-card context. Each transaction is scored on its own. Real fraud detection
uses sequence and account history, such as a burst of small test purchases or a sudden
change in geography, and none of that is available here.

\subsection{Construct Validity}

The cost model is assumed, not measured. The EUR 3 review cost in
Section~\ref{sec:cost} is a placeholder. The conclusion that the deployed threshold is
within 0.3 percent of cost-optimal depends on it, and would change with a realistic figure.

Amount is also an imperfect proxy for fraud loss. Charging a missed fraud its full
transaction amount ignores recovery, chargeback mechanics, and reputational cost.

\section{Practical Recommendations}

\begin{enumerate}
  \item Prefer class weighting over any resampling method on data this imbalanced.
        It matched or beat every alternative, adds no hyperparameters, and introduces no
        synthetic data.
  \item Choose the deployed model on engineering criteria, since the statistical
        criteria do not separate them. Inference speed, calibration quality, memory
        footprint, and model size are all reasonable tie-breakers, and by those measures
        LightGBM remains a good choice.
  \item Collect more fraud labels before further tuning. The learning curve is
        still rising at 100 percent of the data, so extra labelled fraud would give more
        than extra algorithmic effort.
  \item Set the threshold from a real cost model, and revisit it as review capacity
        changes. The framework in Section~\ref{sec:cost} needs only the true review cost
        substituted in.
  \item Report intervals, not single numbers. On a dataset with fewer than 500
        positives, a bare PR-AUC figure looks more precise than it really is.
\end{enumerate}
